% !PRINTED 358
\weberpage{358}
%<<CONT
lineare Function, die drei willkürliche Constanten enthält, so
folgt aus (1):
\[
f = x_1 (V + 2pu + x_1p^2) - (u + px_1)^2,
\]
was wieder von der Form (1) ist.

Ist nun $y_1 = 0$ die Gleichung einer zweiten von $x_1$ verschiedenen Doppeltangente, so können wir die Constanten in $p$ (und zwar noch auf unendlich viele verschiedene Arten) so wählen, dass der Kegelschnitt $u + px_1 = 0$ durch die beiden Berührungspunkte von $y_1$ geht, und wir können daher annehmen, dass schon in der Form (1) die Function $u$ so gewählt sei. Dann muss in denselben Punkten auch die cubische Form $V$ verschwinden.

Aber noch mehr: Da die Linie $y_1 = 0$ Doppeltangente sein soll, so muss, wenn wir $x_1, y_1$ und irgend eine dritte davon unabhängige lineare Function $z$ als Variable einführen, in den Berührungspunkten
\[
f'(x_1) = 0, \quad f'(z) = 0
\]
sein, und daraus folgt nach (1):
\[
V'(x_1) = 0, \quad V'(z) = 0,
\]
d. h. die Curve dritter Ordnung $V = 0$ wird von der Linie $y_1 = 0$ in den Berührungspunkten mit $f$ berührt, und dies ist nur möglich, wenn $V$ zerfällt und die Function $y_1$ als Theiler enthält. Hiernach erhalten wir eine neue Gestalt der Function $f$:
\weq{2}{f = x_1 y_1 v - u^2,}
worin $u, v$ Functionen zweiten Grades sind.

Sind $x_1, y_1$ irgend beliebige lineare, $u, v$ quadratische Formen von drei Variablen, so stellt die durch (2) bestimmte Function $f$, gleich Null gesetzt, eine Curve vierter Ordnung dar, von der $x_1$ und $y_1$ zwei Doppeltangenten sind.

Sind umgekehrt $x_1$ und $y_1$ zwei beliebige Doppeltangenten einer Curve vierter Ordnung $f = 0$, so kann $f$ auf die Form (2) gebracht werden.

Denken wir uns die Co\"efficienten von $x_1$ und $y_1$ als variabel, und lassen diese Grössen so variiren, dass $x_1$ und $y_1$ sich auf der Curve $f$ schneiden, so wird ein Doppelpunkt eintreten.

Die Doppeltangenten arten aus in zwei von dem Doppelpunkte auslaufende Tangenten.
